% !TEX program = xelatex
\documentclass[12pt,a4paper]{ctexart}

\usepackage{fontspec}
\setCJKmainfont{SimSun}[BoldFont=SimHei,ItalicFont=KaiTi]
\setCJKsansfont{SimHei}
\setCJKmonofont{FangSong}

\usepackage[top=28mm,bottom=28mm,left=28mm,right=28mm]{geometry}
\usepackage{setspace}
\usepackage{fancyhdr}
\usepackage{hyperref}
\usepackage{perpage}
\usepackage{xcolor}

\MakePerPage{footnote}

\newcommand{\circledfootnotenumber}{%
  {\fontspec{SimSun}\ifcase\value{footnote}\or ①\or ②\or ③\or ④\or ⑤\or ⑥\or ⑦\or ⑧\or ⑨\or ⑩\else\arabic{footnote}\fi}}
\renewcommand{\thefootnote}{\circledfootnotenumber}

\hypersetup{colorlinks=true, linkcolor=black, urlcolor=blue!50!black}

\setlength{\parindent}{2em}
\setlength{\parskip}{0pt}
\linespread{1.35}

\pagestyle{fancy}
\fancyhf{}
\fancyfoot[C]{\thepage}
\renewcommand{\headrulewidth}{0pt}

\begin{document}

\begin{center}
    {\zihao{2}\bfseries 从《摆脱贫困》看以人民为中心的发展思想\par}
    \vspace{0.5em}
    {\zihao{3}\bfseries ——兼论精准扶贫的基层治理逻辑\par}
    \vspace{1.2em}
    {\zihao{-4}
    姓名：\underline{\makebox[3cm]{纪文龙}} \qquad
    学号：\underline{\makebox[4cm]{2024012842}} \qquad
    班级：\underline{\makebox[3cm]{行健烽火4}}\par}
    \vspace{1.0em}
\end{center}

\begin{abstract}
《摆脱贫困》记录了闽东地区走出贫困的过程，这本书也是理解中国特色社会主义现代化基层治理的生动样本。它并非单纯陈述宏大叙事，而是由群众路线出发，落实到干部作风建设，并进一步延伸到地方产业选择，将价值立场、组织作风与发展路径编织在一起。基于本门课程的教学内容，本文从“以人民为中心”在基层治理中的具体展开入手，分析《摆脱贫困》中的基层治理逻辑。
本文认为：中国的脱贫成就并不在于简单的资源投入，而在于形成了一套系统化的治理机制，能够回应群众真实需求、调动群众主体力量，并在干部与群众之间形成良性互动。
\end{abstract}

\noindent \textbf{关键词：}《摆脱贫困》；以人民为中心；精准扶贫；基层治理；共同富裕

\vspace{1.5em}

\section*{一、问题的提出：以《摆脱贫困》观察基层治理}

在课程学习的过程中，我逐渐产生了一个问题：“以人民为中心”这样的宏大命题，如何在现实中转化为可操作、可衡量的治理方法？在传统的体制及其对应的经济学视野中，政策往往偏好总量数据和增长效率。然而，如果我们将视野拉回真实的闽东大山，拉回中西部深度贫困地区，人们需要的并不是这些数据，而是更切实的生存诉求：谁的生活条件能够得到实质性提升？谁获得了受教育的机会？谁的意见能够被倾听？

教材中指出，人民立场是习近平新时代中国特色社会主义思想的出发点。\footnote{本书编写组：《习近平新时代中国特色社会主义思想概论（2023年版）》，高等教育出版社、人民出版社，2023年，第4章“坚持以人民为中心”。}这告诉我们，评估任何一项公共政策的成效，必须建立在微观个体的生活福祉和社会层面的客观改善之上，而不能停留在报表中的数字增长上。

正因如此，我选择将《摆脱贫困》这本书作为切入点。它是一份生动的文献，展现了地方领导者在错综复杂的人文环境、恶劣的自然条件下，如何开展具体的管理。在基层实际运作中，贫困绝非只有收入这一个单一维度，它是由地理阻隔、传统观念、基础建设不足等多重因素相互交织而成的结构性网络。面对这种系统性危机，单靠“发钱发物”的传统救济，往往会陷入边际效益递减的陷阱。真正需要做的，是把资源投入从单纯外部输血，转变为基层党组织的重构以及农户个体主观能动性的调动，这样才能在泥泞中开辟出可持续发展的道路。

这就使闽东的脱贫实践成为“人民”价值内核的体现。在这一治理场景中，群众从被动等待行政命令的客体，转变成了政策设计和自我发展的主体。干部的职责也从高高在上的管理者，转变为能够深入田间地头、体察民情、协调关系的服务者。这种治理角色的转变，揭示了新时代基层治理的规律。

\section*{二、为民造福不是空泛口号}

在这本书中，一条清晰的价值主线贯穿始终：如何处理好公权力与人民福祉之间的关系。书中谈到“为官一任，造福一方”时，并不是在进行传统的儒家道德修养，而是阐明了现代公共治理的价值根基——公权力需要服务于人民福祉，需要转化为群众生活条件的改善。\footnote{习近平：《摆脱贫困》，福建人民出版社，1992年。}这种底层逻辑，也呼应了课程中以人民为主体的思想内涵。如果剥离掉政策话语的包装，就会发现“人民”并不是一个空洞的政治概念，而是具体到每一个对饮水、道路、教育和农产品销路有所期盼的人。

具体到地方政府的行政实践，这种“造福”观念包含两个维度。首先是破除形式主义。我们可以看到一些反面的例子，在部分地区，有些干部容易将贫困的根源归因于自然条件恶劣、群众素质低，以此来敷衍应付，或者堆砌台账材料。其次是防范“大包大揽”。假如干部只忙着帮农户盖房子，却没有培育起村民的自主发展动力，也没有帮助他们形成适应市场竞争的能力，一旦资源被撤出，返贫几乎成为必然。学术界在梳理宁德时期的治理模式时，也常把它提炼成一种以人民根本利益为核心的治理模式。\footnote{陈振明、吕志奎：《〈摆脱贫困〉中的地方治理思想研究》，《马克思主义与现实》2017年第1期。}它要求干部把价值关怀内化为日常行动，在制度运作中体现对人的尊重。

这也引发了我们对公共治理成效的思考：我们可能很容易关注“外部输入端”，却常常忽视了群众自我发展的能力。如果一个贫困村庄虽然账面上的收入达标，但是在产业基础、协作组织和基础教育等方面依然是一片荒漠，那么这种脱贫就是无源之水、无本之木，难以为继。由此可见，摆脱贫困与实现共同富裕是两个相互联系、相互促进的过程：前者通过兜底来消除生存层面的赤字，后者通过优化资源配置来提升发展能力，让每个人都能平等地参与现代化、分享现代化的成果。

\section*{三、精准扶贫的治理逻辑：把“人民”具体到每一户、每一项需求}

我国扶贫政策从“粗放式灌溉”走向“精准式滴灌”，是一场重大的治理范式转型。传统区域性扶贫开发虽能在宏观上改善基础设施，但对于隐蔽在山褶子里的真正贫困个体，往往存在“溢出效应有限”和“靶向偏离”的缺陷。汪三贵等学者的研究明确指出，精准扶贫在基层落地时，核心考验在于精准识别、扶持与考核这三重维度，这极度依赖地方治理机制的自我创新。\footnote{汪三贵、郭子豪：《论中国的精准扶贫》，《贵州社会科学》2015年第5期。}这与我们概论课中强调的“实事求是”、“问题导向”方法论完全咬合。如果无法建立起一套穿透基层的真实信息网络，政策资源就很容易被层层耗散在“平均用力”的数字游戏中，甚至造成扶贫资源的错配与浪费。

精准扶贫之所以能够将“以人民为中心”的抽象信念转化为制度性现实，在于它解构了宏观的“平均数”。它把政策焦点对准了具体的村庄、家庭乃至每一个活生生的人。不同的农户面临着不同的致贫根源——有些是因为家庭主要劳动力遭遇重大疾病，有些是因为身处偏远山区缺少稳定的农业产业，有些则是缺乏现代劳动技能。这种高度碎片化、差异化的治理诉求，必须通过网格化的“一户一策”来予以回应。葛志军等学者在针对西部两村的田野调查中也发现，在这一精密行政运作的背后，往往存在着贫困户参与权虚化、政策刚性有余而柔性不足、以及基层工作强度超载等多重微观困境。\footnote{葛志军、邢成举：《精准扶贫：内涵、实践困境及其原因阐释——基于宁夏银川两个村庄的调查》，《贵州社会科学》2015年第5期。}这些实证材料恰恰警示我们，单凭顶层设计的完美蓝图是不够的，基层执行中的调适能力和对当地实情的反馈机制才是政策成败的胜负手。

同时，精准扶贫并不是没有问题的理想过程。邓维杰在讨论精准扶贫难点时提到，实践中可能出现规模排斥、区域排斥和识别排斥等问题。\footnote{邓维杰：《精准扶贫的难点、对策与路径选择》，《农村经济》2014年第6期。}这一观察极为深刻：以人民为中心从来都不是一个一蹴而就的终极状态，而是一个不断发现偏差、校正偏轨的动态纠偏过程。如果底层的建档立卡程序脱离了村民委员会的民主监督与日常公示，沦为干部关门评定的“密室决策”，那么所谓的“精准”也只是数据层面的自娱自乐。越是面临纷繁复杂的利益分配，就越要依靠公开透明的信息披露与群众的实质参与。

简而言之，精准扶贫的实质就是用“求真务实”的底层逻辑去重塑基层官僚系统的行为模式。它以深度调研取代主观臆测，以分类治理取代简单的一刀切，以农户的真实获得感作为考核的最高判准。它的生命力在于将国家行政资源的宏观调配与乡村社会的微观生态进行无缝缝合，而不是在无休止的“填表迎检”中消耗治理效能。

\section*{四、调查研究与群众路线：基层治理不能只在材料里完成}

这就必然引申出另一个关于治理工具的重要命题：调查研究。宁德工作期间所留下的“四下基层”（信访接待下基层、现场办公下基层、调查研究下基层、政策宣传下基层）实践，已被历史证明是打破官僚体制层级阻隔、重塑党群联系的有效制度安排。\footnote{申坤：《新时代背景下习近平调查研究重要论述的价值意义》，《山东干部函授大学学报》2020年第3期。}这不仅是政治修辞，它代表了一种深刻的认识论：群众路线绝不是浮光掠影的“走基层”，更不是应付差事的文牍工作，而是干部必须具备穿透行政迷雾、在复杂现场厘清真问题的反思能力。

一旦失去了高频度的深度调研，基层政权在运转中极易滋生两种异化倾向：要么是上级决策层对底层社会运作规律进行悬空式的逻辑推演，导致政策南橘北枳；要么是基层部门将应付各种量化硬性指标视为唯一的理性选择，陷入材料扶贫、PPT展示的空转漩涡。这在复杂的贫困治理中代价极大。一户人家到底面临什么难处，仅凭表册上的几项收支数据常常会被遮蔽；一个村社的协同能力与产业前景，也绝非几份规划书所能描绘。这必须要求公职人员走进人情世故、走进地头民心，在田野现场的直接对话中发掘那些被行政语言过滤掉的“地方性知识”。

这正是《摆脱贫困》跨越时代并持续散发活力的原因。书中没有以居高临下的拯救者姿态去打量闽东这块落后之地，而是不厌其烦地反思地方治理主体的思想局限、干群关系的日常磨合以及本土资源的合理配置。在我国宣告消除绝对贫困并全面转向“乡村振兴”的当下，这一逻辑更显重要。绝对贫困的消解并不意味着基层社会治理的毕其功于一役，相反，防止规模性返贫、乡村公共秩序的重组以及产业内生循环的构建，均需要更为细密、长效的制度安排。

\section*{五、从脱贫攻坚到共同富裕：发展成果怎样更稳定地留在人民生活中}

从历史宏观维度来看，脱贫攻坚不仅是中国摆脱贫穷的面貌改变，更是百年来中国共产党对于社会正义和人的解放的承诺兑现。杨明伟在其思想史的回顾中阐明，消除贫困并朝向共同富裕，是中国人民近代追求复兴的底层动力，是极具现代性张力的社会变革。\footnote{杨明伟：《百年奋斗史中的摆脱贫困迈向共同富裕》，《红旗文稿》2021年第6期。}然而，置于“中国式现代化”的宏观视域下，这只是一曲壮丽序幕的终了。共同富裕被赋予了超越传统资本主义福利国家的全新社会理想，它不仅指涉国民财富总量的增加，更要求我们在区域、城乡以及阶层之间搭建起能力提升的普惠通道，消弭机会的不平等。

我们应当警惕将共同富裕简单化为均贫富的“福利民粹主义”，亦或停留在分发式的慈善救济上，它是以高质量经济循环和公共资源合理配置为前置条件的长期渐进过程。中西部许多刚刚摘帽的脆弱地区，其产业结构单一、公共财政薄弱、教育资源依然存在显著的代际差距。胡蝶等人的研究也强调，只有深化教育等社会资源的精准配置，阻断贫困在代际之间的刚性传递，才能赋予弱势群体自我发展的长效动能。\footnote{胡蝶：《扶贫必扶智：教育精准扶贫是摆脱贫困的内生动力》，《改革与开放》2016年第21期。}这一点至关重要——真正稳定的脱贫并非只看当年收入是否越过了某条贫困线，而是要看其下一代能否凭借公平的教育和健康保障，顺利跨入更高维度的社会生产链条。

如果我们把《摆脱贫困》、精准扶贫、乡村振兴与共同富裕放在同一个大历史视角下审视，会清晰地展现出一条“从生存权保障到发展权共享”的演进脉络。其首要目标是满足生理底线的安全感，随后过渡到通过教育赋权与产业导入实现能力的横向扩容，最终指向以高质量分配结构达成现代化成果的全社会共享。习近平总书记在面对突发公共卫生事件与经济波动等复杂局面时，曾坚决部署必须坚守如期打赢脱贫攻坚战的底线目标，这深刻揭示了我们决策层将社会正义与大众基本民生置于最高权衡序列的执政追求。\footnote{习近平：《在决战决胜脱贫攻坚座谈会上的讲话》（2020年3月6日），《求是》2020年第6期。}这种部署体现了决策层将人民生活水平与社会正义置于最优先地位的执政理念。

作为青年学生，我们在反思上述治理图景时，应努力超越教材概念的空洞拼接，而将目光聚焦于那些真实发生的社会关系变迁。比如：在我们的家乡，回乡创业的年轻人是否获得了合理的政策孵化？乡村小学在推行“合并分流”时是否伤害了偏远村庄家庭的就学便利？中西部乡镇的年轻干部在面对考核压力时如何保持为民服务的初心？这些具体的基层图景，才正是评判“以人民为中心”是否在基层土地扎根生长的终极考题。

\section*{结语}

《摆脱贫困》中虽载有诸多特定时空下的政策细节，但其核心提供的实际上是理解当代中国国家治理如何建构行政合法性与政策执行力的线索。它提醒我们，以人民为中心的发展思想并非自我标榜的政治修辞，而是深深嵌入在干部的日常办公逻辑、政策的靶向纠偏、长期的调查研究以及群众的主体动员过程中的制度设计。脱贫攻坚的大规模实践，正是将“人民”这一虚化、同质化的代名词，具体还原为面临各异生存困境的千万户个体。

没有任何一项制度设计能具有天然的完美性，任何宏大的改革措施也必将在微观执行中遭遇到地方性摩擦与官僚异化的挑战。新时代人民立场的应有之义，在于决策者和执行者能够始终维持虚怀若谷的纠偏能力，不回避实操中的漏洞，敢于在直面真实问题的碰撞中优化制度。由《摆脱贫困》所开启，经精准扶贫而丰富，并朝向乡村振兴及共同富裕蔓延的历史进程，在逻辑深处是一条绵延不断的探索之路：如何以党的全面领导保证政治方向，以人民的主体地位提供治理热能，以实事求是作为工具箱，让现代化建设的每一滴甘霖，都能毫无阻碍地滋润进普通人的日常柴米油盐之中。

\end{document}
